\documentclass[conference]{IEEEtran}

% Packages
\usepackage{amsmath}
\usepackage{amssymb}
\usepackage{graphicx}
\usepackage{booktabs}
\usepackage{array}
\usepackage{url}
\usepackage{cite}



% 
\title{Underwater Vessel Noise Classification Using the ShipsEar Dataset:\\
       Background, Acoustic Theory, and Formal Problem Definition}

\author{
  \IEEEauthorblockN{Author Name(s)}
  \IEEEauthorblockA{
    \textit{Department / Lab Name}\\
    \textit{Institution Name}\\
    City, Country\\
    email@institution.edu
  }
}

% 
\usepackage[hidelinks]{hyperref}
\begin{document}
\maketitle

% 
\begin{abstract}
This report establishes the scientific and engineering context for underwater
vessel noise classification using the ShipsEar dataset~\cite{shipsear2016}.
We review the dataset's recording protocol and hydrophone specifications,
examine the physical mechanisms governing underwater acoustic propagation
(sound speed, transmission loss, multipath interference, propeller cavitation,
machinery harmonics, and blade-rate modulation), characterise the five vessel
categories in terms of their distinct spectral signatures, and formally define
the classification task as a supervised multi-class pattern recognition
problem. This background forms the theoretical foundation for all subsequent
feature extraction, model training, and evaluation stages of the project.
\end{abstract}

\begin{IEEEkeywords}
underwater acoustics, vessel noise classification, ShipsEar, propeller
cavitation, LOFAR, DEMON, hydrophone, multi-class classification
\end{IEEEkeywords}

% 
\section{Introduction}

Identifying a vessel from its acoustic signature --- without visual contact ---
is a long-standing problem in naval sonar, maritime surveillance, and marine
environmental monitoring. A submerged hydrophone records a composite signal
arising from propeller cavitation, engine vibration, hull flow noise, and
propeller blade-rate modulation. Each vessel type imprints a characteristic
spectral fingerprint on this signal, making automatic classification feasible
through signal processing and machine learning.

The ShipsEar dataset~\cite{shipsear2016}, introduced by Santos-Dom\'inguez
\textit{et al.}\ in 2016, provides the first publicly available benchmark
corpus for this task, covering five operationally distinct vessel classes
recorded in a realistic estuarine environment. This part of the project
reviews the dataset, the relevant acoustic physics, and states the
classification objective precisely in mathematical terms.

% 
\section{Dataset Overview: ShipsEar}

\subsection{Recording Setup}

ShipsEar was recorded in the R\'ia de Vigo estuary, northwestern Spain.
An omnidirectional HTI-96-MIN hydrophone was deployed at approximately
1.5\,m depth. The sensor exhibits a flat frequency response from 2\,Hz
to 30\,kHz with a sensitivity of $-$64\,dB re 1\,V/$\mu$Pa, making it
well-suited for capturing the full spectral bandwidth of marine vessel noise.
All 90 recordings are stored as uncompressed 16-bit PCM WAV files sampled at
$f_s = 52{,}000$\,Hz. Recording durations vary from approximately 15\,s to
over 3\,min, reflecting real-world operational diversity.

\subsection{Vessel Taxonomy}

The dataset is partitioned into five mutually exclusive vessel classes,
summarised in Table~\ref{tab:classes}. The dataset is intentionally
class-imbalanced: Class\,A (motorboats) is overrepresented relative to
Classes\,C and E, an important consideration for classifier design and
evaluation.

\begin{table}[ht]
  \caption{ShipsEar Vessel Classes and Acoustic Characteristics}
  \label{tab:classes}
  \centering
  \setlength{\tabcolsep}{4pt}
  \begin{tabular}{>{\bfseries}c l p{4.6cm}}
    \toprule
    \textbf{Label} & \textbf{Vessel Type} & \textbf{Acoustic Characteristics} \\
    \midrule
    A & Motorboats
      & High-frequency broadband noise; intense propeller cavitation at high RPM; non-stationary. \\[4pt]
    B & Mussel Boats
      & Prominent diesel LOFAR tonal lines; slow blade-rate harmonics; stable narrowband peaks. \\[4pt]
    C & Fishing Vessels
      & Mixed broadband and tonal content; winch and gear noise; variable engine load. \\[4pt]
    D & Passengers / Ferries
      & Strong, stable engine harmonics (50--500\,Hz); consistent blade-rate due to fixed routes. \\[4pt]
    E & Ocean Liners / Tugboats
      & Very low fundamental frequency; highest source level; complex multi-shaft harmonics. \\
    \bottomrule
  \end{tabular}
\end{table}

% 
\section{Underwater Acoustic Propagation}

\subsection{Speed of Sound}

Sound travels in seawater at approximately 1{,}500\,m/s --- roughly 4.4 times
faster than in air (343\,m/s). The empirical Mackenzie
equation~\cite{mackenzie1981} gives the sound speed $c$ as a function of
temperature $T$ ($^\circ$C), salinity $S$ (ppt), and depth $z$ (m):

\begin{align}
  c \approx\; & 1448.96 + 4.591T - 0.05304T^2 + 0.0002374T^3 \nonumber \\
              & + 1.340(S-35) + 0.0163z \quad \text{[m/s].}
  \label{eq:mackenzie}
\end{align}

The higher propagation speed relative to air results in longer acoustic
wavelengths at any given frequency, enabling low-frequency vessel harmonics
to propagate over kilometre-scale ranges with relatively low attenuation.

\subsection{Transmission Loss}

In shallow coastal waters such as the R\'ia de Vigo, acoustic energy
undergoes cylindrical spreading supplemented by frequency-dependent
absorption~\cite{urick1983}:

\begin{equation}
  \mathrm{TL} = 10 \log_{10}(r) + \alpha r \quad \text{[dB],}
  \label{eq:tl}
\end{equation}

where $r$ is the source-to-receiver range in metres and $\alpha$ is the
absorption coefficient in dB/m. Because $\alpha$ increases with frequency,
high-frequency components (above $\sim$10\,kHz) are strongly attenuated at
modest ranges, concentrating discriminative vessel information below 10\,kHz.
This motivates resampling the ShipsEar data from 52\,kHz to a lower target
rate (e.g., 22{,}050\,Hz) without information loss.

\subsection{Multipath Interference}

In shallow estuarine environments, acoustic energy undergoes multiple
reflections from the pressure-release sea surface and the impedance-mismatch
seabed. These multipath arrivals produce constructive and destructive
interference at the hydrophone, causing frequency-selective fading and
temporal smearing of tonal spectral features. Multipath effects are a
principal source of within-class recording variability in ShipsEar and
motivate the use of segment-level averaging in feature extraction.

% 
\section{Vessel Noise Generation Mechanisms}

Three physical mechanisms dominate underwater vessel radiated noise and
directly motivate the feature extraction choices described in Part\,4.

\subsection{Propeller Cavitation}

Cavitation occurs when the local pressure at a propeller blade surface falls
below the vapour pressure of seawater, causing dissolved gas pockets to form
and subsequently collapse. The implosions produce broadband high-frequency
noise dominant above approximately 1\,kHz. Cavitation intensity scales with
propeller tip speed and is particularly pronounced in small, fast vessels
(Class\,A motorboats). Mel-frequency cepstral coefficients (MFCCs) and
log-mel spectrograms are effective representations for capturing this
broadband energy distribution.

\subsection{Machinery Noise and LOFAR Lines}

Diesel and petrol engines generate narrowband tonal lines at integer multiples
of the shaft rotation frequency. For a shaft spinning at $\omega$ revolutions
per minute, the fundamental machinery frequency is:

\begin{equation}
  f_0 = \frac{\omega}{60} \quad \text{[Hz].}
  \label{eq:lofar_fundamental}
\end{equation}

The resulting harmonic series $\{f_0, 2f_0, 3f_0, \ldots\}$ appears as sharp
peaks in a high-resolution power spectral density (PSD) estimate, commonly
computed via Welch's method. These peaks are referred to as \textit{LOFAR
lines} (Low-Frequency Analysis and Recording) and are among the most
class-discriminative features in ShipsEar, particularly for Classes\,B and D.

\subsection{Blade-Rate Harmonics and DEMON Spectrum}

As each propeller blade passes through the non-uniform wake behind a shaft
strut, it generates a periodic pressure pulse. For a propeller with $N_b$
blades on a shaft rotating at $\omega$\,RPM, the fundamental blade-rate
frequency is:

\begin{equation}
  f_\mathrm{BR} = N_b \cdot \frac{\omega}{60} \quad \text{[Hz].}
  \label{eq:blade_rate}
\end{equation}

Blade-rate harmonics modulate the amplitude envelope of the broadband
cavitation noise. The \textit{DEMON} (Detection of Envelope Modulation on
Noise) technique demodulates this envelope via a Hilbert transform, then
computes the PSD of the envelope signal to reveal $f_\mathrm{BR}$ and its
harmonics. DEMON features are highly discriminative for propeller-driven
vessels and complement LOFAR analysis.

% 
\section{Acoustic Profiles of Vessel Classes}

Building intuition for class-specific spectral characteristics is essential
before modelling decisions are made.

\textbf{Class\,A --- Motorboats:} Smallest and fastest vessels. High
rotational speed produces intense cavitation; dominant spectral energy is
broadband above 2\,kHz. Signals are non-stationary due to variable throttle.

\textbf{Class\,B --- Mussel Boats:} Medium-sized working vessels at slow,
steady speed. Diesel LOFAR lines are sharp and stable. Blade-rate frequency
is low ($\approx$1--4\,Hz). The combination of strong narrowband peaks and
low broadband floor makes this class relatively easy to separate.

\textbf{Class\,C --- Fishing Vessels:} Spectral non-stationarity arises from
switching between transit and active trawling modes. Winch and mechanical gear
noise introduce additional tonal components absent in other classes.

\textbf{Class\,D --- Passengers/Ferries:} Large vessels with fixed regular
routes and consistent engine loads. Stable harmonic series dominates
50--500\,Hz. May be confused with Class\,B due to overlapping blade-rate
range; separation relies on source level and harmonic spacing differences.

\textbf{Class\,E --- Ocean Liners/Tugboats:} Highest source levels in the
dataset. Very low fundamental frequencies result from large, slow-turning
propellers. Multi-shaft configurations produce complex overlapping harmonic
patterns. Cavitation is present at full operational load.

% 
\section{Formal Problem Definition}

\subsection{Input Space}

Let $\mathbf{x} \in \mathbb{R}^{N}$ denote a fixed-length audio segment
captured by the hydrophone, where:

\begin{equation}
  N = T_s \times f_s,
  \label{eq:input_dim}
\end{equation}

$T_s$ is the segment duration in seconds (e.g., 3\,s) and $f_s$ is the
(resampled) sampling rate. Each segment is obtained by applying a sliding
window with hop size $H < T_s$ to a raw recording, discarding any terminal
remainder shorter than $T_s$.

\subsection{Label Space and Classification Objective}

Each segment is assigned a discrete class label $y$ from the finite set:

\begin{equation}
  y \in \mathcal{Y} = \{\,A,\; B,\; C,\; D,\; E\,\}.
  \label{eq:label_set}
\end{equation}

The learning objective is to find an optimal classifier $f^*$ that minimises
the expected classification loss $\mathcal{L}$ over the joint data
distribution $p(\mathbf{x}, y)$:

\begin{equation}
  f^* = \operatorname*{arg\,min}_{f}\;
        \mathbb{E}_{(\mathbf{x},\,y)\sim p}\!\left[\,
        \mathcal{L}\!\left(f(\mathbf{x}),\, y\right)\right].
  \label{eq:objective}
\end{equation}

In practice, $f$ operates on a compact feature representation
$\boldsymbol{\phi}(\mathbf{x}) \in \mathbb{R}^{D}$ with $D \ll N$, extracted
via handcrafted signal processing (MFCCs, LOFAR, DEMON) or a learned
convolutional encoder (log-mel spectrogram).

\subsection{Key Assumptions and Constraints}

The following assumptions govern the problem formulation:

\begin{enumerate}
  \item \textbf{Single-channel input:} Each segment is captured by one
        omnidirectional hydrophone; no spatial beamforming is assumed.

  \item \textbf{Mutual exclusivity:} Labels are exhaustive and mutually
        exclusive --- each segment belongs to exactly one class in
        $\mathcal{Y}$.

  \item \textbf{Recording-level splitting:} Train, validation, and test sets
        are partitioned at the \emph{recording} level (not the segment level)
        to prevent data leakage between overlapping windows derived from the
        same source file.

  \item \textbf{Class imbalance:} ShipsEar is imbalanced; model training uses
        class-weighted loss, and evaluation relies on macro-averaged F1-score
        and AUC-ROC rather than raw accuracy alone.

  \item \textbf{Reproducibility:} All random operations (splitting,
        weight initialisation) are seeded with fixed values from
        $\{0, 7, 21, 42, 99\}$; mean $\pm$ standard deviation is reported
        across five runs.
\end{enumerate}

\subsection{Performance Metrics}

Let $\hat{y} = f(\mathbf{x})$ denote the predicted label. The primary
evaluation metrics are:

\begin{align}
  \text{Accuracy} &= \frac{1}{|\mathcal{T}|}\sum_{i\in\mathcal{T}}
                     \mathbf{1}[\hat{y}_i = y_i], \label{eq:acc} \\[4pt]
  \text{Macro-F1} &= \frac{1}{|\mathcal{Y}|}\sum_{c\in\mathcal{Y}}
                     \frac{2\,P_c\,R_c}{P_c + R_c}, \label{eq:f1} \\[4pt]
  \text{Macro-AUC} &= \frac{1}{|\mathcal{Y}|}\sum_{c\in\mathcal{Y}}
                      \mathrm{AUC}_{\mathrm{OVR}}(c), \label{eq:auc}
\end{align}

where $P_c$ and $R_c$ are the precision and recall for class $c$,
$\mathcal{T}$ is the held-out test set, and AUC$_\mathrm{OVR}$ is computed
under a one-vs-rest scheme.

% 
\section{Conclusion}

This part has established the complete scientific context required before any
data processing or modelling. The ShipsEar dataset provides 90 hydrophone
recordings across five acoustically distinct vessel classes in a realistic
shallow-water environment. The dominant noise mechanisms --- propeller
cavitation, LOFAR engine harmonics, and blade-rate modulation --- motivate the
feature engineering choices pursued in Parts\,3--4. The formal problem
definition introduces the mathematical notation, loss framework, evaluation
metrics, and experimental constraints used throughout the remainder of the
project.

% 
\begin{thebibliography}{9}

\bibitem{shipsear2016}
D.~Santos-Dom\'inguez, S.~Torres-Guijarro, A.~Cardenal-L\'opez, and A.~Pena,
``ShipsEar: An underwater vessel noise database,''
\textit{Applied Acoustics}, vol.~113, pp.~64--69, 2016.
\url{https://doi.org/10.1016/j.apacoust.2016.06.008}

\bibitem{urick1983}
R.~J. Urick,
\textit{Principles of Underwater Sound}, 3rd~ed.
New York, NY, USA: McGraw-Hill, 1983.

\bibitem{ross1976}
D.~Ross,
\textit{Mechanics of Underwater Noise}.
New York, NY, USA: Pergamon Press, 1976.

\bibitem{mackenzie1981}
K.~V. Mackenzie,
``Nine-term equation for sound speed in the oceans,''
\textit{Journal of the Acoustical Society of America},
vol.~70, no.~3, pp.~807--812, 1981.

\end{thebibliography}

\end{document}
